\documentclass[11pt]{article}
\usepackage[a4paper,margin=25mm]{geometry}
\usepackage[T1]{fontenc}
\usepackage{lmodern,microtype,amsmath,amssymb,amsthm,mathtools}
\usepackage{booktabs,longtable,array,xcolor,enumitem,fancyhdr,multicol}
\usepackage[breaklinks,colorlinks=true,linkcolor=blue!45!black,urlcolor=blue!45!black,citecolor=blue!45!black]{hyperref}
\usepackage{listings}
\hypersetup{pdftitle={Sixteen integer unknowns for an explicit consistency equation},pdfauthor={Research synthesis for the ProveIt project}}
\lstset{basicstyle=\ttfamily\small,breaklines=true,columns=fullflexible,keepspaces=true,frame=single,rulecolor=\color{black!15}}
\setlength{\emergencystretch}{2em}
\setlist{itemsep=3pt,topsep=5pt}
\newtheorem{theorem}{Theorem}[section]
\newtheorem{lemma}[theorem]{Lemma}
\newtheorem{proposition}[theorem]{Proposition}
\newtheorem{corollary}[theorem]{Corollary}
\theoremstyle{remark}\newtheorem{remark}[theorem]{Remark}
\newcommand{\Z}{\mathbb Z}\newcommand{\N}{\mathbb N}
\newcommand{\Pos}{\mathbb Z_{>0}}\newcommand{\Q}{\mathbb Q}
\newcommand{\Sq}{\operatorname{Sq}(\Z)}
\newcommand{\ZFC}{\mathsf{ZFC}}\newcommand{\Con}{\operatorname{Con}(\ZFC)}
\newcommand{\SignedOps}{1590}
\newcommand{\pin}{\texttt{8c6eec9d19ef80538c8e61ff6e0180040583446a}}
\pagestyle{fancy}\fancyhf{}\fancyhead[L]{\small Explicit Diophantine consistency forms}\fancyhead[R]{\small ProveIt synthesis}\fancyfoot[C]{\thepage}
\setlength{\headheight}{14pt}
\title{Sixteen integer unknowns for an explicit consistency equation\\[5pt]\large and simpler polynomial presentations of $\operatorname{Con}(\mathsf{ZFC})$}
\author{Research synthesis for the ProveIt project\\\small Prepared with AI assistance; ordinary mathematical proofs and reproducible artifacts}
\date{22 October 2026}
\begin{document}
\maketitle
\begin{abstract}
We combine the current ProveIt constructions with a positivity-normalization
lemma for Pell equations. The resulting explicitly specified integer polynomial
$P_{16}$ has sixteen unrestricted integer unknowns and satisfies
\[
 \ZFC\vdash\bigl(\nexists\mathbf x\in\Z^{16}:P_{16}(\mathbf x)=0\bigr)
 \longleftrightarrow\Con.
\]
Its degree is exactly $406528\cdot5^{32}+10008$ and its delivered arithmetic
circuit uses \SignedOps{} binary operations. The reduction removes a guard
witness from the existing seventeen-integer construction: a signed Pell
coordinate can be replaced by an arbitrarily large positive coordinate in
the same residue class. A parity-sensitive radical weight lowers the degree
further. We also obtain a 61-operation system of twenty-two quadratics in
31 positive unknowns, with a 126-operation quartic sum of squares, and an
identical 108-operation evaluator for the existing degree-1081 polynomial
in twelve positive unknowns. A separate complete native coefficient tuple
reduces the radix from the previous best $2^{1,953,125}$ to $2^{1,012,625}$.
The arithmetic circuits, finite coefficient
recipes, exact algebra checks, source pins, and reproduction scripts are
included. These are improvements to explicit presentations in the inspected
repository, not global minimality claims or a proof of consistency. The new
proofs are ordinary mathematics; the inherited native compiler and Jones
representation theorems remain identified dependencies.
\end{abstract}
\section*{Contents}
\setcounter{tocdepth}{1}
\begin{multicols}{2}\small
\makeatletter\@starttoc{toc}\makeatother
\end{multicols}
\newpage

\section{What is improved, and what ``simplest'' means}
The reviewed repository revision is \pin. The target directory is
\texttt{con-zfc-diophantine-forms}, under
\path{SetTheory/Cardinals/docs/reports/hilbert-tenth-problem}; its audited
README describes seven wave-4 packages, numbered 24--30 \cite{repo}.
The source snapshot is retained in the archive. Different presentations
optimize different quantities; number of variables, degree, coefficient
height, and arithmetic evaluation cost cannot be collapsed into one ordering.

\begin{center}\small
\begin{tabular}{p{43mm}p{39mm}p{55mm}}
\toprule
Measure & Inspected predecessor & This synthesis\\
\midrule
Explicit unrestricted-integer construction & 17 variables; 20,607 operations & 16 variables; \SignedOps{} operations\\
Degree of that construction & $424296\cdot5^{33}+53824$ & $406528\cdot5^{32}+10008$\\
31-positive-variable quadratic system & 22 equations; 65 operations & 22 equations; 61 operations\\
31-positive-variable quartic & 130 operations; 306 symbolic monomials & 126 operations; 259 symbolic monomials\\
12-positive-variable degree-1081 polynomial & 111 operations & 108 operations; exactly the same polynomial\\
Associated three-equation system & 107 operations & 104, or 103 after a justified coordinate shift\\
Native coefficient radix & $2^{1,953,125}$ & $2^{1,012,625}$; separate complete tuple\\
\bottomrule
\end{tabular}\end{center}

All operation counts charge binary addition, subtraction, and multiplication,
including scalar multiplication and squaring. Sharing an already computed
value and comparing two system sides are free. Fixed integer coefficients
are leaves, not additional unknowns. Fixed exponentiation is expanded into
multiplications. Thus these are arithmetic-circuit sizes, not bit-complexity
bounds and not measurements of the work needed to construct the coefficients.
The monomial comparison keeps the four native coefficient ports symbolic.

The repository's nine-natural-variable constructions remain better in
variable count over $\N$; its 67-operation single equation remains better
in operation count, and its 59-operation systems remain better in that
metric. We do not claim a smaller universal variable bound. In particular,
Sun's 2026 preprint gives an effective general representation theorem in
ten unrestricted integers \cite{sun}; that different compiler is not
instantiated here. Our sixteen-variable claim concerns the concrete,
fully bound constructions supplied with this report.

We use $\N=\{0,1,2,\ldots\}$ and $\Pos=\{1,2,\ldots\}$.
An equation over $\Pos$ must not silently be read over $\Z$. Our main
$P_{16}$ does allow every one of its sixteen variables to range over all
of $\Z$.

\section{Fixed source instance and logical dependencies}\label{sec:binding}
The logical statement $\Con$ means that no finite proof of contradiction
exists in the fixed recursive presentation of ZFC used by the repository.
The calibrated machine is the literal 215-state machine
\texttt{candidate-k215/k215.tm}. Its Lean endpoint states
\begin{lstlisting}
blank_halts_iff_zfc_refutation :
  BlankHalts machine215 <-> BProv ZFCax_s [] Form.fBot
\end{lstlisting}
The source file is included. This report neither rebuilds that development
nor relies on a metatheoretic assertion that Lean is sound to infer a theorem
of ZFC. We use the ordinary arithmetic calibration mirrored by the endpoint,
then compose ordinary mathematical reductions formalizable in ZFC.

The native compiler contributes four fixed integers
\[
 \beta=B_0-1,\qquad K_0=D_C+B_0D_R,\qquad M_C,\qquad
 M_F=M_F^{\rm native}+B_0-1,\qquad B_0=2^{5^{10}}.
\]
Throughout the main construction these are exactly the complete compact943
tuple of package 30. In particular $M_F$ is the \emph{shifted} mask. The
four gzip payloads are included, together with the original binding, layout,
ordered 943-tile list, and coefficient proofs. The raw sizes are
\begin{center}\small
\begin{tabular}{lrr}
\toprule Integer & Binary digits & Raw bytes\\\midrule
$\beta$ & 9,765,625 & 1,220,704\\
$K_0$ & 13,201,251 & 1,650,157\\
$M_C$ & 9,765,625 & 1,220,704\\
$M_F$ & 9,765,626 & 1,220,704\\\bottomrule
\end{tabular}\end{center}
The total is 5,312,269 raw bytes, stored in 20,260 gzip bytes. Byte hashes
are checked during reproduction. Compression does not change coefficient
height. No coefficient is chosen by testing whether ZFC is consistent.

The imported compiler theorem identifies positive solutions of the parent
quartic with blank halting of K215. Its hypotheses include the native-window
language, no-carry bounds, synchronization, rank control, and the Pell
arithmetization; they are not assertions about arbitrary choices of four
integer ports. The source proofs are bundled, and the new transformations
preserve this one coherent tuple. The Jones reduction is the concrete
quartic-to-ratio construction of \cite{jones}, mirrored by
\texttt{Jones1982.Nine.ratio\_sound},
\texttt{ratio\_complete}, and \texttt{wset\_iff\_nine} in the repository.

\section{A cheaper positive quartic}\label{sec:quartic}
All thirty-one variables in this section are strictly positive:
\[
\begin{gathered}
 J,F,\alpha,z,f,h,i,o,j,s,w,\tau,\eta,\zeta,y,\gamma,\\
 N,Y,P,c,a,X,D,C_2,L,F_2,V,g,I,B_{\rm aux},\Delta.
\end{gathered}
\]
Use only the affine abbreviations
\[
 q=\beta J+1,\quad k=\eta+\zeta,\quad H=4a+3,\quad
 C=q-F-\alpha,\quad E=a-Y.
\]
The new system consists of the following twenty-two quadratic equations:
\begin{alignat}{2}
 L&=q^2-1,&\qquad N&=qL,\notag\\
 X&=wq,& Y&=s(N+q),\notag\\
 E&=XY,&c&=kY+\eta,\notag\\
 P&=E(c-\eta),&\tau^2&=P(P+k)+1,\notag\\
 D&=X+ac+\gamma H,&\Delta&=a^2+H,\notag\\
 C_2&=c^2,&D^2&=\Delta C_2+1,\notag\\
 k&=jc-V+1+hE,&&\notag\\
 jc-V&=(q-F)N-CL+(M_C+qM_F)J,&&\notag\\
 (K_0+w)C&=F+z(q-1),&&\notag\\
 F_2&=f^2,&V&=of-c,\notag\\
 g&=y^2-V^2,&I&=iC_2,\notag\\
 B_{\rm aux}&=I^2,&B_{\rm aux}&=\Delta(F_2-1),\notag\\
 y^2-B_{\rm aux}g&=1.&&\label{eq:quartic-system}
\end{alignat}
Let $r_1,\ldots,r_{22}$ be these left sides minus right sides, in row order,
and define $Q_{31}^{*}=\sum r_j^2$. The JSON gives this order unambiguously.

\begin{proposition}
The positive zero set of \eqref{eq:quartic-system} is in bijection with the
positive zero set of package 30's parent quartic. Consequently
\[
 \ZFC\vdash(\nexists\mathbf x\in\Pos^{31}:Q_{31}^{*}(\mathbf x)=0)
 \longleftrightarrow\Con.
\]
The degree is exactly four. The system has a 61-operation circuit, and
the displayed sum of squares has a 126-operation circuit.
\end{proposition}
\begin{proof}
Write a subscript $o$ for parent coordinates. Apply the simultaneous affine
change
\[
 N_o=N+q,\qquad E_o=a-Y,\qquad P_o=P+k,\qquad z_o=z+1,
\]
leaving all other variables fixed. Every parent residual becomes exactly
the corresponding residual of \eqref{eq:quartic-system}. All twenty-two
identities are checked as sparse integer-polynomial identities with symbolic
coefficient ports. The inverse change is
$N=N_o-q$, $Y=a-E_o$, $P=P_o-k$, $z=z_o-1$.

A positive new zero restores $E_o=XY>0$; the other three restored
coordinates are visibly positive. Conversely the parent equations give
$N=qL>0$, $Y=sN_o>0$, and $P=E_okY>0$. For the fourth coordinate we
use the parent's established transport-quotient lemma, also used in its
retained59 construction. After the Pell/compiler reconstruction it gives
$C>0$, $X=wq=2^R$, $R>3q$, and $F<q$. Thus $w>F$, and
\[
 0<(K_0+w)C-F=(z_o-1)(q-1)
\]
forces $z_o\ge2$. This is the necessary positive-domain argument; positivity
of $z_o$ alone would not suffice.

Each of the four changes saves one gate in the scheduled system evaluation.
For example $P_o=E_o(c-\eta)+k$ loses its final addition, while its norm
replaces $P_o-k$ by $P+k$ at the same cost. The $N$ change deletes two
operations and adds the one operation in $N+q$. The result is
$61=33$ multiplications $+28$ additions/subtractions. Forming twenty-two
residuals, twenty-two squares, and their sum costs $22+22+21=65$ more.
Finally $\tau$ appears only as $\tau^2$ in $r_8$; the coefficient of
$\tau^4$ in the sum of squares is exactly one.
\end{proof}

\section{A finite, fully specified Jones index}\label{sec:index}
In this and later sections all letters are newly scoped. Put $\nu=31$ and
$\mathcal L=5^{32}$. Define
\[
 Q(x,n_1,\ldots,n_{31})=
 Q_{31}^{*}(n_1+1,\ldots,n_{31}+1),
\]
independently of the dummy input $x$. At the all-zero natural tuple the fifth
residual is $-1$, so $Q(x,0)>0$. Thus $Q$ satisfies the normalization
hypothesis of the concrete Jones theorem, and has degree four.

For $\alpha\in\N^{32}$ with $|\alpha|\le4$, let $p_\alpha$ be its
ordinary coefficient and put
\[
 c_\alpha=(4-|\alpha|)!\prod_{j=0}^{31}\alpha_j!\,p_\alpha,
 \qquad e(\alpha)=\sum_{j=0}^{31}\alpha_j5^j.
\]
There are 391 nonzero symbolic records, each an explicitly listed polynomial
in $(\beta,K_0,M_C,M_F)$; the shifted quartic has 571 terms with these
ports symbolic. The new coefficient file is \texttt{data/coefficients31.json}.
It is freshly derived from $Q_{31}^{*}$, not obtained by reusing the parent's
old coefficient list after changing the parent.

For a record $c_\alpha=\sum_{\mathbf e} a_{\mathbf e}
\prod_{j=1}^{4}b_j^{e_j}$, use the authentic port bit lengths $\ell_j$
in Section~\ref{sec:binding} and the finite integer bound
\[
 |c_\alpha|\le H_\alpha:=
 \sum_{\mathbf e}|a_{\mathbf e}|2^{\sum_j e_j\ell_j}.
\]
The supplied exact arithmetic certificate verifies
$H_\alpha+1<2^{45933757}$ for every record. Unlisted coefficients are zero.
Set, once and for all,
\begin{align*}
 z_*&=2^{45933757},\\
 e_*(T)&=\sum_{|\alpha|\le4}(z_*+c_\alpha)
                T^{\mathcal L-e(\alpha)},\\
 y_*&=e_*(2z_*),\qquad
 u_* =\sum_{j=1}^{31}(2z_*)^{5^j}.
\end{align*}
All exponents are fixed nonnegative integers; in particular
$e(\alpha)\le4\cdot5^{31}<\mathcal L$.
The sum has $\binom{36}{4}=58905$ terms. Its base-$2z_*$ digits are
strictly between 1 and $2z_*$, in distinct positions. These definitions
satisfy precisely the admissible index hypothesis for the normalized
quartic $Q$. The unused $u_*$ is specified to make that interface explicit.

This is a finite definition of literal integer coefficients. Materializing
the enormous integer $y_*$ is unnecessary to specify the polynomial. It is
not an unspecified recursive-enumeration index, an existential parameter,
or a value defined by whether a computation halts.

\section{The eight-coordinate ratio predicate}\label{sec:ratio}
Use the eight base coordinates
\[
 (\epsilon,g,h,s,w,\phi,i,j).
\]
Initially, Jones's predicate uses $\epsilon,\phi\ge0$ and
$g,h,s,w,i,j\ge1$. Later $h,i,j$ will range over all integers.
All other symbols in this section are polynomial abbreviations:
\begin{gather*}
 b=y_*+1+\epsilon,\quad B=b^5,\quad q=B^{\mathcal L},\quad
 \ell=\sum_{r=1}^{31}B^{5^r},\quad
 \lambda=\sum_{r=0}^{4\mathcal L-1}B^r,\\
 e=e_*(B),\quad \theta=B-2z_*,\quad N=q^{16},\quad
 G_c=32bB^{5^{31}}.
\end{gather*}
Define $R$ by
\begin{align*}
 R={}&\bigl[g+eq^3+\ell q^5+
  \{2(e-z_*\lambda)(1+B+g)^4+\lambda B+\lambda Bq^4\}q^7\bigr](N^2-N)\\
 &+\bigl[q^3-b\ell+\ell+\theta\lambda q^3+(B-2)q^8\bigr](N^2-1).
\end{align*}
The remaining abbreviations are
\begin{align*}
 U&=N^2w,&Y&=N^2s,&M&=RY,\\
 P&=2M^2U,&K&=R+1+h(P-1),&A&=M(U+1),\\
 B_1&=2R+1,&W_1&=bw,&C&=B_1+W_1+\phi,\\
 J&=4A-5,&Q_1&=3W_1C-2(W_1^2-1),&D&=(A^2-1)C^2+1,\\
 E&=iJDC^2,&F&=(A^2-1)E^2+1,&G_1&=A+F(CD+1-A),\\
 H&=B_1+jC,&I&=(G_1^2-1)H^2+1.&&
\end{align*}
Finally put
\begin{gather*}
 a_0=(P^2-1)K^2+1,\qquad a_1=DFI,\\
 m=K^2G_c-4(C-KY)^2G_c-8gK^2,\\
 D_0=JF,\qquad T_0=Q_1F+J(H-C).
\end{gather*}
The imported concrete Jones equivalence says
\begin{equation}
 \begin{gathered}
 \exists\text{ base coordinates in the initial domain}:\quad
 a_0,a_1\in\Sq,\ D_0\mid T_0,\ m>0\\
 \Longleftrightarrow\quad \exists\mathbf n\in\N^{31}:Q(1,\mathbf n)=0.
 \end{gathered}
 \tag{JR}\label{eq:JR}
\end{equation}
This is the ratio predicate underlying the nine-variable theorem, not a
claim that an arbitrary nine-variable polynomial has this structure.

\section{Removing three positivity restrictions}\label{sec:normalization}
\begin{lemma}\label{lem:denom}
For every integer assignment of the eight base coordinates, $D_0\ne0$.
Replacing a nonzero $i$ by $|i|$ preserves $a_0,a_1,D_0,T_0,m$.
\end{lemma}
\begin{proof}
$J=4A-5\ne0$. If $F=0$, then $(1-A^2)E^2=1$, forcing $A=0$ and
$E=\pm1$. But then $J=-5$ and $E=iJDC^2$ is divisible by 5, a
contradiction. The only nonsquared occurrence of $i$ is in $E$, which
enters the five listed quantities only through $E^2$.
\end{proof}

\begin{lemma}\label{lem:h}
Suppose $\epsilon,\phi\ge0$, $g,s,w\ge1$, $a_0\in\Sq$, and $m>0$.
Then $h\ge1$. In particular the proof does not assume a sign for $j$.
\end{lemma}
\begin{proof}
The index gives $b\ge2$, $B=b^5>2z_*$, and $e>0$.
Since $G_c>0$ and $g>0$, the margin implies $K\ne0$,
\[
 g<G_c/8=4bB^{5^{31}},\qquad 4(C-KY)^2<K^2.
\]
Writing $t=5^{31}$, we have
$1+B+g<1+B+4B^{t+1}\le5B^{t+1}<B^{t+2}\le B^{5t}=q$.
Also $\ell<B^{t+1}$, hence $b\ell<B^{t+2}\le q$.
In the left bracket for $R$, the possibly negative term
$-2z_*\lambda(1+B+g)^4$ is dominated by $\lambda Bq^4$.
Every remaining term there is nonnegative and $g$ is positive. The
right bracket is positive as well. Thus $R\ge N^2-N\ge2$.

Now $Y,U\ge1$, $C=2R+1+bw+\phi>0$, and
$P=2R^2Y^2U\ge2R^2>R+2$. The margin excludes $K\le0$:
if $K<0$, then $|C-KY|> |K|$, and $K=0$ was already excluded.
If $h\le0$, then $1\le K\le R+1<P$. The case $K=1$ would give
$P-1\mid R$, impossible because $0<R<P-1$. For $2\le K<P$,
\[
 (PK-1)^2<(P^2-1)K^2+1<(PK)^2,
\]
as the two differences are $K(2P-K)>0$ and $K^2-1>0$.
This contradicts squarehood of $a_0$.
\end{proof}

The next lemma supplies the new witness saving.
\begin{lemma}[Positive Pell coordinates in a prescribed residue class]\label{lem:pellpositive}
Let $a\ge2$ be an integer and $d=a^2-1$. If
$x^2-dy^2=1$ with $x>0$ and $y\in\Z$, then for every modulus $M\ge1$
there are solutions $(x',y')$ with $x',y'>0$, arbitrarily large $y'$, and
$x'\equiv x$, $y'\equiv y\pmod M$.
\end{lemma}
\begin{proof}
The element $u=a+\sqrt d$ is a unit of norm one in $\Z[\sqrt d]$ and
its multiplication matrix has determinant one. In the finite group
$\mathrm{SL}_2(\Z/M\Z)$ some positive power $u^r$ is the identity.
Multiply $x+y\sqrt d$ by $u^{rn}$. Its coordinates are integers, its norm
stays one, and both coordinates retain their residues modulo $M$.
The real number $x+y\sqrt d$ is positive: its conjugate has the same
sign and their sum is $2x>0$. Since $u>1$, the second coordinate is
\[
 \frac{(x+y\sqrt d)u^{rn}-(x+y\sqrt d)^{-1}u^{-rn}}{2\sqrt d}
 \longrightarrow+\infty.
\]
The first coordinate is positive as well. This proves the claim, including
negative or zero starting second coordinates.
\end{proof}

\begin{proposition}[The $j$ guard is unnecessary]\label{prop:j}
The existential predicate \eqref{eq:JR} is unchanged if its base domain is
replaced by
\[
 \epsilon,\phi\ge0,\qquad g,s,w\ge1,\qquad
 h,j\in\Z,\qquad i\in\Z\setminus\{0\}.
\]
\end{proposition}
\begin{proof}
Only soundness of the larger domain needs proof. Lemmas~\ref{lem:denom}
and \ref{lem:h} normalize $i$ and show $h\ge1$. We now have
$A>1$, $C>B_1>0$, and $G_1>1$.
The congruences $F\equiv1\pmod D$ and $F\equiv1\pmod J$ follow
from $E=iJDC^2$. Consequently $D_0\mid T_0$ implies
$F\mid H-C$. Modulo $D$ we have $G_1\equiv1$, hence $I\equiv1$.
Modulo $F$, $G_1\equiv A$ and $H\equiv C$, hence $I\equiv D$.
Therefore $D,F,I$ are pairwise coprime positive integers. Squarehood of
$DFI$ forces each of them to be a square.

Choose the positive integer $x=\sqrt I$; then
$x^2-(G_1^2-1)H^2=1$. Apply Lemma~\ref{lem:pellpositive} with
$a=G_1$ and modulus $CF$. Choose $H'>B_1$ congruent to $H$ modulo $CF$,
and set $j'=(H'-B_1)/C$. This is a positive integer, since
$H\equiv B_1\pmod C$. The new $I'$ is a square, so $DFI'$ is a square.
Also
\[
 T_0'-T_0=J(H'-H)\quad\hbox{is divisible by }JF.
\]
All other ratio conditions remain unchanged. We have recovered a tuple in
the original domain, to which \eqref{eq:JR} applies. The map is existential:
it changes the numerical value of $I$, while preserving its required
squarehood and the required congruences.
\end{proof}

\section{Encoding the remaining guards}\label{sec:guards}
Introduce six unrestricted integer variables $z_1,\ldots,z_6$ and put
\[
 (d_1,d_2,d_3,d_4,d_5,d_6)=
 (3\epsilon+2,3g-1,3s-1,3w-1,3\phi+2,3m-1).
\]
Starting with $C_0=1$, define
\[
 B_r=1+4d_rC_{r-1}^2z_r^2,\qquad
 C_r=C_{r-1}B_r\quad(1\le r\le6),\qquad C_*=C_6.
\]
\begin{lemma}\label{lem:stack}
Subject to $z_1\cdots z_6\ne0$, the condition $C_*\in\Sq$ is equivalent,
existentially in the six $z_r$, to
$\epsilon,\phi\ge0$, $g,s,w\ge1$, and $m>0$.
Moreover $C_*\equiv1\pmod4$ for every integer assignment, without any
squarehood or nonvanishing hypothesis.
\end{lemma}
\begin{proof}
Each $B_r\equiv1\pmod4$ is nonzero, and
$\gcd(B_r,C_{r-1})=1$. Thus the $B_r$ are pairwise coprime.
If their product is a square, each $|B_r|$ is a square; a negative square
is 0 or 3 modulo 4, so each $B_r$ is a positive square. Because
$C_{r-1}z_r\ne0$, this forces $d_r\ge0$. Its residue 2 modulo 3 excludes
zero, giving exactly the six desired inequalities.

Conversely each positive $d_r\equiv2\pmod3$ is nonsquare, so its Pell
equation $u^2-d_rv^2=1$ has a nontrivial positive solution. A power of its
unit is the identity modulo $2|C_{r-1}|$, by the same finite-group argument
as above. We may therefore take $v$ a positive multiple of
$2|C_{r-1}|$ and set $z_r=v/(2C_{r-1})\ne0$. Then $B_r=u^2$.
Induction proves the converse. The final congruence follows directly from
the formulas.
\end{proof}

\section{One divisor for divisibility and nonvanishing}\label{sec:merger}
The existing signed construction uses the degree-five intersective polynomial
\[
 \Pi(t)=(t^3-t-1)(t^2+23).
\]
For completeness the necessary arithmetic is proved here.
\begin{lemma}\label{lem:intersective}
$\Pi$ has no rational root but has a root modulo every nonzero integer.
\end{lemma}
\begin{proof}
The monic cubic has only $\pm1$ as possible rational roots, and neither
works. The quadratic has no real root. The cubic discriminant is $-23$.
For an odd prime $p\ne23$, a cubic root modulo $p$ is simple and lifts.
If there is no cubic root, the cubic is irreducible. Frobenius permutes its
three roots by a 3-cycle, so their Vandermonde product is fixed; its square
is $-23$. The quadratic consequently has a simple root and lifts instead.
At $p=23$, the cubic root $3$ is simple because its derivative is $26$.
At $p=2$, use $-23\equiv1\pmod8$. A root modulo $2^e$ for $e\ge3$
lifts to $2^{e+1}$ by choosing one of $t$ and $t+2^{e-1}$.
Roots modulo arbitrary integers now follow by the Chinese remainder theorem.
\end{proof}

Define the polynomial
\[
 \mathcal R(D,T,v)=
 \bigl((T+Dv)^3-D^2(T+Dv)-D^3\bigr)
 \bigl((T+Dv)^2+23D^2\bigr).
\]
\begin{lemma}\label{lem:merge}
For integers $D,T,Z$ with $D\ne0$,
\[
 (D\mid T\ \wedge\ Z\ne0)
 \quad\Longleftrightarrow\quad
 \exists v\in\Z:\ ZD^5\mid\mathcal R(D,T,v).
\]
Also $\mathcal R(D,T,v)\ne0$ for every integer $v$.
\end{lemma}
\begin{proof}
The identity $\mathcal R=D^5\Pi(T/D+v)$ proves nonvanishing.
If $T=Dq$ and $Z\ne0$, choose $v$ by
Lemma~\ref{lem:intersective} so that $Z\mid\Pi(q+v)$.
Conversely $Z=0$ would force $\mathcal R=0$. If $D\nmid T$, choose
a prime with $\alpha=v_p(D)>\beta=v_p(T)$. The first term in each
factor of $\mathcal R$ has uniquely least valuation: respectively
$3\beta$ and $2\beta$. Thus $v_p(\mathcal R)=5\beta<5\alpha$,
contradicting divisibility by $D^5$.
\end{proof}

\section{The polynomial in sixteen unrestricted integers}\label{sec:main}
The sixteen variables are exactly
\[
 \epsilon,g,h,s,w,\phi,i,j,z_1,z_2,z_3,z_4,z_5,z_6,v,u\in\Z.
\]
Using Sections~\ref{sec:ratio}--\ref{sec:guards}, put
\begin{gather*}
 Z=i z_1z_2z_3z_4z_5z_6,\qquad S=ZD_0^5,\qquad
 T=\mathcal R(D_0,T_0,v),\qquad X=Su-T,\\
 w_1=C_*+1,\qquad w_2=w_1(P^2+K^2+2),\\
 A_0=S^2C_*,\qquad A_1=(Sw_1)^2a_0,\qquad A_2=(Sw_2)^2a_1,\\
 V=X^2+A_2-A_0-A_1,\qquad
 W=V^2+4X^2A_2-4A_0A_1.
\end{gather*}
Our equation is the ordinary integer-polynomial equation
\begin{equation}
 \boxed{\ P_{16}=W^2-16V^2X^2A_2=0.\ }\label{eq:p16}
\end{equation}
Every intermediate letter is an abbreviation. In particular none adds a
quantifier. The coefficient leaves are completely fixed in
Sections~\ref{sec:binding} and \ref{sec:index}.

\begin{lemma}[Parity-sensitive separation]\label{lem:weights}
At every integer assignment,
\[
 |w_1|>\sqrt{|C_*|},\qquad
 |w_2|>\sqrt{|C_*|}+|w_1|\sqrt{|a_0|}.
\]
\end{lemma}
\begin{proof}
Because $C_*\equiv1\pmod4$, either $C_*\ge1$ or $C_*\le-3$.
In both cases $(C_*+1)^2>|C_*|$. For arbitrary real $P,K$,
\[
 \sqrt{|(P^2-1)K^2+1|}\le P^2+K^2+1.
\]
Indeed bound the absolute radicand by $(P^2+1)K^2+1$ and square the
nonnegative right side; the remaining difference is
$P^4+K^4+P^2K^2+2P^2+K^2\ge0$. Multiplying by $|w_1|$ proves
the second strict inequality. Negative values of $w_1,w_2$ cause no issue.
\end{proof}

\begin{lemma}[Three-root finalizer]\label{lem:finalizer}
For fixed earlier coordinates, \eqref{eq:p16} has an integer solution in $u$
if and only if
\[
 S\ne0,\qquad S\mid T,\qquad C_*,a_0,a_1\in\Sq.
\]
\end{lemma}
\begin{proof}
Lemmas~\ref{lem:denom} and \ref{lem:merge} give $T\ne0$. If $S=0$,
then $P_{16}=T^8\ne0$. Assume $S\ne0$, and choose complex roots
$r_0^2=C_*$, $r_1^2=a_0$, $r_2^2=a_1$. Exact multiplication gives
\[
 P_{16}=\prod_{\delta_0,\delta_1,\delta_2\in\{\pm1\}}
 [X-S(\delta_0r_0+\delta_1w_1r_1+\delta_2w_2r_2)].
\]
A zero makes one signed radical sum equal to the rational number $u-T/S$.
Two distinct sign patterns on nonzero roots cannot give the same sum:
take their largest differing index; a nonzero square root of an integer has
absolute value at least one, and Lemma~\ref{lem:weights} makes its weighted
contribution exceed the sum of all preceding contributions. The triangle
inequality proves this even for negative radicands.

The field generated by the three roots is a finite Galois extension of
$\Q$, and its automorphisms change their signs. Every automorphism fixes
the rational sum; separation therefore forces it to fix each root individually.
Each root lies in $\Q$, hence in $\Z$ because its square is an integer.
Zero roots are already rational. Thus all three radicands are nonnegative
integer squares and $T/S$ is an integer. Conversely, choose their integer
roots and set $u=T/S+r_0+w_1r_1+w_2r_2$.
\end{proof}

\begin{theorem}[Main equivalence]\label{thm:main}
For the fixed coefficients defined above,
\[
 \ZFC\vdash
 \bigl(\nexists\mathbf x\in\Z^{16}:P_{16}(\mathbf x)=0\bigr)
 \longleftrightarrow\Con.
\]
\end{theorem}
\begin{proof}
A zero gives the conditions of Lemma~\ref{lem:finalizer}.
Since $D_0\ne0$, $S\ne0$ implies $i z_1\cdots z_6\ne0$.
Lemma~\ref{lem:merge} recovers $D_0\mid T_0$, and
Lemma~\ref{lem:stack} recovers $\epsilon,\phi\ge0$, $g,s,w\ge1$, $m>0$.
Proposition~\ref{prop:j} restores $h,i,j>0$ existentially. The concrete
Jones equivalence, the quartic equivalence, and the K215 calibration now
give a finite ZFC refutation.

Conversely a refutation supplies a positive quartic solution and then a
base tuple satisfying \eqref{eq:JR}. Lemma~\ref{lem:stack} supplies six
nonzero guard witnesses with $C_*\in\Sq$. Lemma~\ref{lem:merge} supplies
$v$ with $S\mid T$; here $S\ne0$. Lemma~\ref{lem:finalizer} supplies $u$.
Negating the existence equivalence proves the claim. The steps use integer
arithmetic, finite-group arguments, elementary Pell theory, and finite Galois
theory, all formalizable in ZFC.
\end{proof}

\section{Exact degree and circuit specification}\label{sec:degree}
Write $\mathcal L=5^{32}$. Degrees below treat every fixed coefficient as
degree zero. The relevant source degrees are
\begin{center}\small
\begin{tabular}{lrlr}\toprule
Quantity & Degree & Quantity & Degree\\\midrule
$R$ & $235\mathcal L$ & $P$ & $950\mathcal L+3$\\
$K$ & $950\mathcal L+4$ & $a_0$ & $3800\mathcal L+14$\\
$a_1$ & $24640\mathcal L+68$ & $m$ & $2221\mathcal L+11$\\
$D_0$ & $6875\mathcal L+20$ & $C_5$ & $363$\\
$C_*$ & $2221\mathcal L+1102$ & $w_1$ & $2221\mathcal L+1102$\\
$w_2$ & $4121\mathcal L+1110$ & $S$ & $34375\mathcal L+107$\\
$T$ & $34375\mathcal L+105$ & $X$ & $34375\mathcal L+108$\\
$A_2$ & $101632\mathcal L+2502$ & $P_{16}$ & $406528\mathcal L+10008$\\\bottomrule
\end{tabular}\end{center}
The first five stacks obey $c_r=3c_{r-1}+3$, $c_0=0$; the final stack
uses the margin. $A_2$ has strictly higher degree than $X^2,A_0,A_1$.
Thus the leading terms of $V,W,P_{16}$ are respectively $A_2,A_2^2,A_2^4$.
There is no top-degree cancellation. An independent diagonal coefficient
certificate modulo $1,000,003$ gives the nonzero final leading residue
$580506$, independent of every fixed coefficient leaf. Therefore
\[
 \deg P_{16}=9\,465\,217\,590\,332\,031\,250\,000\,010\,008.
\]

The delivered circuit is \texttt{data/signed16.json}. Its selected evaluator
uses \SignedOps{} live binary gates. The large digit polynomial is evaluated
using the complete homogeneous polynomial $h_4$. Set
$t_j=B^{5^j}$, $v_j=B^{5^{31}-5^j}$, and $t=t_{31}$.
Then
\[
 e_*(B)=t\left[z_*h_4(t,v_0,\ldots,v_{31})+
 \sum_{|\alpha|\le4}c_\alpha t^{4-|\alpha|}\prod_jv_j^{\alpha_j}\right].
\]
The exponent of a summand is exactly
$5\cdot5^{31}-\sum_j\alpha_j5^j$.
Initialize $h_0=1$ and $h_1=h_2=h_3=h_4=0$.
The ascending recurrence $h_r\leftarrow h_r+xh_{r-1}$, for $r=1,2,3,4$,
evaluates the first term; multivariate Horner arithmetic evaluates the sparse
correction. Shared fifth powers give $t_j$, while
$v_j=t_j^4v_{j+1}$. Finally
\[
 \lambda=\prod_{j=0}^{31}(1+t_j+t_j^2+t_j^3+t_j^4)(1+q)(1+q^2).
\]
These identities are polynomial identities; the output circuit has no division.

\section{An identical twelve-variable polynomial in 108 operations}\label{sec:positive12}
This improvement preserves the exact fixed polynomial of package 30.
Its variables are $J,F,\alpha,z,h,i,T,s,w,\eta,\zeta,t\in\Pos$.
The four constants are the same complete tuple as before. Define
\begin{gather*}
 q=\beta J+1,\ C_w=q-F-\alpha,\ X=wq,\ Y=sq^3,\ E=XY,\\
 k=\eta+\zeta,\ c=kY+\eta,\ a=E+Y,\ H=4a+3,\ \Delta=a^2+H,\\
 \ell=EkY,\quad A=1+\ell(\ell+k),\quad G=1+\Delta c^2,\quad b_0=X+ac,\\
 R=((q-1)(q-F)+\alpha)(q^2-1)+(M_C+qM_F)J,\\
 N_k=k-hE-R,\qquad N_t=(K_0+w)C_w+(q-F)-z(q-1),\\
 I=ic^2,\quad D=1+\Delta I^2,\quad Q_0=\Delta(D-1),\quad d=Q_0-1,\\
 v=cT,\quad b=c+RD,\quad M=k(I^2+T^2),\quad L=HM+1,\quad W=b_0+Ht.
\end{gather*}
A complete scalar evaluator for the norm $\mathcal N$ is
\begin{align*}
 H_2&=H^2,& H_Q&=H_2Q_0,& H_P&=H_QDv^2,&H_U&=H_Qb^2,\\
 n&=(HM)^2A,&d_n&=dn,&H_e&=H_P+H_U-H_2,&j_0&=4H_PH_U,\\
 W_2&=W^2,& L_2G&=L^2G,&\mathrm{Tr}&=2d(W_2+L_2G),&&\\
 K_r&=d(W_2-L_2G),&K_n&=K_r^2,&A_A&=2(d_n+H_e),&&\\
 C_c&=(d_n-H_e)^2-j_0,&a_t&=\mathrm{Tr}-A_A,&b_t&=C_c-K_n,&&\\
 c_f&=16d_nj_0,&v_r&=a_t(a_t\mathrm{Tr}+2b_t)-c_f,&&&&\\
 u_r&=b_t^2-a_t^2K_n.&&&&&&
\end{align*}
The final norm is
\[
 \mathcal N=u_r(u_r+\mathrm{Tr}v_r)+K_nv_r^2.
\]
The equation is
\[
 \boxed{\ P_{12}=N_kN_t(1+\mathcal N)-1=0.\ }
\]
It has degree 1081 and uses 108 operations. The system
$N_k=1$, $N_t=1$, $\mathcal N=0$ uses 104 operations, with maximum
residual degree 1072. Replacing the transport variable by $z_o=z+1$ gives
the comparison $(K_0+w)C_w=F+z(q-1)$ and saves one more operation.
The positive-domain justification is the transport argument already given.

\begin{proposition}
The 108-operation circuit is identical as an integer polynomial to the
111-operation source circuit. Its consistency equivalence is unchanged.
\end{proposition}
\begin{proof}
There are precisely three algebraic savings. First share $H^2Q_0$ before
forming $H_P,H_U$, saving one multiplication. Next compute
$[d(W^2-L^2G)]^2$ instead of $d^2(W^2-L^2G)^2$, saving one multiplication.
Finally Horner-evaluate
$u_r^2+\mathrm{Tr}u_rv_r+K_nv_r^2$ as in the display, saving one gate.

For an independent norm derivation, let
$P_a=Q_0Dv^2$, $U_a=Q_0b^2$, $e_a=P_a+U_a-1$, and $B_a=2Q_0vb$.
Then $\mathcal N$ is the eight-factor product
\[
 \prod_{u^2=G,\ \tau^2=A,\ f^2=D}
 \bigl[d(W-Lu-HM\tau)^2-H^2e_a+H^2B_af\bigr],
\]
where each root has both signs. Eliminating $\tau,f$ first gives a quartic
in $x=d(W\pm L\sqrt G)^2$ of the form
$(x^2-A_Ax+C_c)^2-c_fx$. The two $x$ values have sum $\mathrm{Tr}$
and product $K_n$. Reducing modulo $x^2-\mathrm{Tr}x+K_n$ gives
$u_r+v_rx$, whose norm is the displayed $\mathcal N$.
All identities hold over polynomial rings, including at repeated or zero roots.
The source's full positive-domain soundness proof is preserved and applies
to the identical polynomial. No new compiler binding is required.
\end{proof}

% Optional report section. Requires amsmath and amssymb; no custom macros.
\section{A smaller complete coefficient tuple}
\label{sec:smaller-coefficients}

The arithmetic constructions and the coefficient compiler can be improved
independently.  The present coefficient construction retains package~27's
ordered Wang943 tile language and its full native/Pell interface, but reduces
the base from $2^{1953125}$ to
\[
 b=125,\qquad L=8101,\qquad d=bL=1012625,\qquad B=2^d.
\]
Thus the base exponent decreases by $48.1536\%$.  The four actual integer
ports occupy $506644$ bytes in minimal unsigned binary form, supplied in
$10502$ bytes of deterministic gzip files.  The largest port has $1015251$
bits.  These are coefficient improvements; no variable, degree, or operation
claim follows merely from changing the tuple.  The main signed-sixteen
instance in this report retains its original fixed tuple.

\subsection{Finite recipe}

Put $\rho=2^{125}$ and $W_j=1024\cdot4^j$ for $0\le j<44$.  Use the
same $943$ ordered tiles $x_{it}$, with ports $t=0,1,2,3$ denoting west,
north, east, south.  The ten-bit code of an edge label $x$ is $x+512$.
The four copy exponent maps are
\[
 e_0(j)=1024+2j,\quad e_1(j)=1004+2j,\quad
 e_2(j)=984+2j,\quad e_3(j)=964+2j \quad(0\le j<10).
\]
Let $E$ consist of selectors $0,\ldots,942$, these forty copy positions,
the dummies $1065,1066$, and anchors $1110,1112,1111,2225$, in that
clause-weight order.  In particular $|E|=989$.  Define
\[
 c_i=1+\sum_{t=0}^3\sum_{j=0}^9
       \operatorname{bit}_j(x_{it}+512)W_{10t+j}
       +\sum_{j=40}^{43}W_j
       \quad(0\le i<943).
\]
At a copy position $e_t(j)$ set $c_{e_t(j)}=W_{10t+j}$; at the four
anchors assign $W_{40},W_{41},W_{42},W_{43}$ respectively; at both
dummies assign zero.  Finally put
\begin{align*}
 \mu&=1022+\sum_{j=0}^{43}W_j,\\
 D_C&=\rho+\rho^{40}+\rho^{61}+\rho^{1113}
       +\sum_{e\in E}c_e(\rho^{3341-e}+\rho^{4459-e}),\\
 D_R&=\rho^{21},\\
 M_F^{\rm nat}&=4+\mu(\rho^{3341}+\rho^{4459})
       +\sum_{j=0}^{9}(\rho^{1004+2j}+\rho^{1045+2j})\\
       &\hspace{14mm}+\rho^{1111}+\rho^{2225}
       +\sum_{z=6685}^{7544}\rho^z,\\
 M_C&=B-1-\sum_{\substack{e\in E\\ e\ne1}}\rho^e-2\rho^{1065}.
\end{align*}
The complete coefficient tuple is
\[
 \boxed{\bigl(B-1,\ D_C+BD_R,\ M_C,\ M_F^{\rm nat}+B-1\bigr).}
\]
The fourth port is already shifted.  All four ports must be replaced
together.  Every coefficient is fixed by bounded finite arithmetic,
independently of the existence of a ZFC refutation.

\subsection{Decoding and geometry}

The mixed-radix clause decoder is unchanged.  Before decoding, the low ten
bits count at most $943$ selected tiles; all other contributions are multiples
of $1024$.  Mask $1022$ forces that count to be zero or one.  After this
extraction each radix-four digit is a sum of at most two Boolean bits, and
its masked low bit forces the required equality.  Thus an occupied cell
contains one tile, its exact forty edge bits, and four anchors equal to one.

The new geometry is summarized by the complete native-basis identities below.
Each row specifies the only nonzero contribution from the indicated source;
an empty entry means no contribution.
\begin{center}
\begin{tabular}{c|c|c|c}
Test target & $D_C$ predecessor & $D_R$ predecessor & Temporal predecessor\\
\hline
$1004+2j$ & $964+2j$ & & $1004+2j$\\
$1045+2j$ & $984+2j$ & $1024+2j$ &\\
$1111$ & $1110$ & & $1111$\\
$2225$ & $1112$ & & $2225$
\end{tabular}
\end{center}
Here $0\le j<10$.  For example, at the north target the unwanted bare
predecessors under shifts $61$ and $1$ occupy the odd intervals
$943,\ldots,961$ and $1003,\ldots,1021$; both miss $E$.  The other rows
follow by the same displayed finite arithmetic.

Write
\[
 S=[0,942]\cup\{964,966,\ldots,1042\}
   \cup\{1065,1066,1110,1111,1112\},\qquad E=S\cup\{2225\}.
\]
The span of $S$ is $1112$.  Clause contamination at a target $t$ would
require $3341-t$ or $4459-t$ in $E-E$.  All these differences exceed
$2225$ except $3341-2225=1116$; that exception is absent because
$2225-1116=1109\notin E$.  At a clause target all bare and spatial
predecessors lie above $E$: the smallest is $3341-1113=2228$.

The band separation $1118$ exceeds the span of $S$, and
$2225-1118=1107\notin E$.  Therefore $1118\notin E-E$.
Since $L>2\max E$, the two possible rotation-hit sets
$(3341-E)\bmod L$ and $(4459-E)\bmod L$ are disjoint.  Every binary
rotation leaves a clean clause band.  Likewise the unique pair of elements
of $E$ with difference $1114$ is $(1111,2225)$, whence
\[
 \{1111-e:e\in E\}\bmod L\ \cap\
 \{2225-e:e\in E\}\bmod L=\{0\}.
\]
The source two-anchor parity argument first excludes nonzero offsets within
a radix digit, including the one possible upper-dummy spill, and this
intersection then forces whole-cell alignment.  The edge and occupancy
conditions give exactly the same Wang matching contract as before.

The conservative source carry bound is also unchanged.  With
\[
 \mathcal M=991\left(2\sum_{e\in E}c_e+6\right)
\]
one has $4\mathcal M+8<2^{125}$ and $2\mu+4<2^{125}$; their relevant
bit lengths are respectively $120$ and at most $99$.  These estimates hold
before selector extraction.  The maximal unshifted convolution exponent
is $4459+2225=6684$.  The $860$ mask fields at $6685,\ldots,7544$ are
therefore identically zero for every aligned intended cell, including every
safe dummy adjustment.  They fit below $L=8101$ and restore the exact
population balance
\[
 \operatorname{pc}(M_F^{\rm nat})=989,\qquad
 \operatorname{pc}(M_C)=d-989,\qquad
 M_C\equiv2\pmod4,\quad M_F^{\rm nat}\equiv4\pmod8.
\]

\subsection{The rank-control extension}

The source chooses a power-of-five cell length to tune the actual packed
rank.  The next argument isolates the necessary arithmetic condition and
allows $L=8101$.

\paragraph{Boolean subset lemma.}
Let $N=5^n$, $M=dN$, and $T=N/5>5d$.  Every residue modulo $M$ is a
Boolean subset sum of the distinct physical weights
\[
 \{B^{4j}:0\le j<T\}\ \cup\ \{B\}.
\]
Indeed $2^{125}\equiv90\pmod{8101}$ and $90^2\equiv-1\pmod{8101}$,
so $B^4\equiv1\pmod{8101}$.  Moreover
\[
 v_5(B^4-1)=v_5(2^{4d}-1)=v_5(2^4-1)+v_5(d)=4.
\]
The Chinese remainder theorem implies that $B^4$ has exact order $T$
modulo $M$.  Its powers are consequently exactly the $T$ residues
$1+5d t$, with $t\bmod T$.

For a target $z$, choose $\epsilon\in\{0,1\}$ so that
$z-\epsilon B\not\equiv0\pmod5$.  Choose $1\le k\le5d$ congruent
to $z-\epsilon B$ modulo $5d$.  Then $k<T$ and $k$ is invertible
modulo $T$.  Set
\[
 u=\frac{z-\epsilon B-k}{5d}\pmod T,\qquad
 r=k^{-1}\left(u-\frac{k(k-1)}2\right)\pmod T.
\]
The $k$ distinct cyclic residues $r,r+1,\ldots,r+k-1$ sum to $u$ modulo
$T$.  The corresponding distinct powers $B^{4j}$, together with the
optional $B$, sum to $z$ modulo $M$.  This uses Boolean choices, with
no repeated physical cell or negative multiplicity.

For the canonical helical witness choose independently sufficiently large
powers of five $w,h$, with $N=wh$ and $w<N/5$.  The selected source
positions are at most $4N/5-4$; their spatial and temporal destinations
remain below $N$.  Thus the inherited exact nonwrapping rank-difference
formula applies.  A dummy change at cell $j$ changes the actual rank by
$-\Gamma B^j$, where
\[
 \Gamma=\rho^{1065}(q^2-1)
        \bigl[1+q(D_C+BD_R+B^w)\bigr],\qquad q=B^N.
\]

To see that $\Gamma$ is a unit modulo $dN$, work modulo $5$ and $8101$.
In both moduli $\rho^2=-1$, and $L,N,w\equiv1\pmod4$, so
$B=q=B^w=\rho$.  The paired clause terms cancel since $1118\equiv2\pmod4$.
The bare terms give $D_C=1+3\rho$ and $D_R=\rho$.  It follows that
\[
 q^2-1=-2,\qquad 1+q(D_C+BD_R+B^w)=-3.
\]
As $\gcd(8101,30)=1$, both factors are units; the factor $\rho^{1065}$
is a unit as well.  The Boolean lemma therefore tunes the actual rank to
$R\equiv dw\pmod{dN}$.  No primality assumption on $8101$ is required.

\subsection{Transfer of the full interface}

The geometry proves soundness by the same clean-band decoder and Wang-plane
argument.  Completeness begins with a genuine finite blank halt, constructs
its padded helical word, and tunes that word's actual rank as above.  Every
dummy change preserves the selector, edge, anchor, mask, and added-zero-field
conditions.  The source positive-slack and half-binomial/Pell constructions
then produce fresh positive witnesses at this actual rank.

For clarity, the extra promise used by signed constructions survives as well.
Put $t=dN$, $q=2^t$, and use the source packed objects
\[
 S_0=(C-1)+qF,\qquad
 T_0=(M_CJ+1)+q(M_F^{\rm nat}J-1),\qquad M_0=q^2.
\]
They satisfy $S_0\mathbin{\&}T_0=0$, $0\le S_0,T_0<M_0$, and
$\operatorname{pc}(T_0)=t+2$.  The exact identity
\[
 R=(M_0-S_0)(M_0-1)+T_0
   =(M_0-1-S_0)M_0+(S_0+T_0)
\]
gives $\operatorname{pc}(R)=3t+2$, while $R\equiv3\pmod4$.
Since $t$ is odd, the source witness $X=2^R$ has square scale
$X/q=2^{R-t}$.  Its ordinary Pell construction supplies the same odd
positive multiplier.  Finally the authentic half-binomial valuation gives
$v_2(Y)=\operatorname{pc}(R)-2=3t$; from $Y=sq^3$ it follows that $s$
is positive and odd.  Thus this tuple supplies the complete source interface,
including the odd-scale promise, and can be used by a compatible arithmetic
skeleton after explicitly rebinding all four ports.

The accompanying programs check all $69230$ native-basis identities,
materialized populations and residues, and an independently reconstructed
coefficient tuple.  They also construct $1029$ Boolean subsets in a small
composite analogue and $32$ subsets for the actual modulus
$d\cdot5^{11}=49444580078125$.  These finite checks support the finite
recipe; the universal mathematical argument is the proof above, composed
with the identified inherited compiler and K215 calibration.  This is
neither a consistency proof nor a new end-to-end proof-assistant certificate.


\section{Verification, reproducibility, and limitations}\label{sec:verification}
The archive is self-contained for its stated finite checks. All release scripts
use Python's standard library; a TeX installation is needed only to rebuild
this PDF. Run from the extracted directory:
\begin{lstlisting}
python3 reproduce.py
pdflatex -interaction=nonstopmode -halt-on-error report.tex
pdflatex -interaction=nonstopmode -halt-on-error report.tex
\end{lstlisting}
The first command rebuilds the circuits and coefficient recipes and runs the
checks. It does not search for a ZFC refutation or a Diophantine solution.

The exact checks include all twenty-two affine residual identities, liveness
and operation counts, an independent expansion of the shifted quartic and all
391 weighted coefficient recipes, the strict index bounds, and the three-root
norm identity as equality of sparse polynomials. The concrete native payloads
are authenticated after decompression. The signed circuit is also compared
with direct displayed formulas at actual-coefficient residues, with literal
58,905-term digit-sum evaluations. Constructed Pell residue examples and
intersectivity checks test transcription at finite values. The twelve-variable
circuit is compared both with the predecessor and with an independent
eight-conjugate evaluator.

Finite checks supplement the universal proofs. They do not establish
consistency, prove a nonhalting computation by testing initial steps, or
replace the imported unbounded native compiler theorem. No new end-to-end
Lean or Rocq certificate is claimed. In particular the paper states a
provable equivalence; it does not assert that ZFC proves either side. Under
the usual consistency hypothesis, the nonexistence assertion is unprovable
in ZFC by the second incompleteness theorem. No stronger two-sided
independence claim is needed here.

The source provenance separates three contributions: the inherited K215,
Pell, compiler, and Jones mathematics; exact refactorings of existing
polynomials; and the new guard-removal and parity-weight arguments.
The supplied sixteen-variable circuit is a concrete new specialization;
we make no assertion of priority for general Pell periodicity, radical norm
merging, or intersective polynomials.

\section{Further research questions}
\begin{enumerate}
\item Can the current eight-coordinate Jones ratio predicate lose another
guard without imposing a restriction on possible bases? In particular,
replacing the base slack by a square is not justified: the recovered base
must also be a power of two, and cofinality of squares alone does not solve
that simultaneous condition.
\item Can a Pell parameter already present in the predicate supply the residue
variable $v$ in Lemma~\ref{lem:merge}? A proof would need residue-class
surjectivity compatible with the full ratio predicate; arbitrary cofinal
Pell solutions are insufficient.
\item Can one remove a register from the 31-variable quadratic system,
giving a normalized quartic in thirty variables? This would lower the Jones
radix exponent by a factor of five, and is more valuable than a small local
gate saving. Positivity of any restored coordinate must be proved.
\item Can the available ten-integer representation theorem be instantiated
with a coefficient recipe and arithmetic DAG for this exact ZFC refutation
predicate? That would address the gap between a general existence theorem
and the concrete sixteen-variable release.
\item Can the fixed-coefficient choice be optimized jointly with the quartic
index rather than measuring the native radix alone? The final $y_*$ size
depends on the shifted quartic coefficient bounds and on $\nu$.
\item Can the positivity-normalization theorem and the parity-sensitive
three-root finalizer be formalized as reusable Lean or Rocq lemmas, then
composed with the existing literal Jones and K215 endpoints? This would
reduce the gap between the new mathematical proof and an end-to-end
machine-checked theorem.
\item Can the native compiler receive a standalone kernel-checked certificate
covering its synchronization, carry bounds, and rank-control argument?
Coefficient hashes and finite geometry checks alone do not certify that
unbounded theorem.
\item Is a single equation below 67 operations attainable under the same
binary-gate convention, or a smaller low-degree system without increasing
the witness count? The present results do not settle either question.
\end{enumerate}

\begin{thebibliography}{9}
\bibitem{repo} ProveIt contributors, \emph{Con(ZFC) in Diophantine form},
wave-4 collection, revision \pin, inspected 22 October 2026.
\href{https://github.com/VladimirReshetnikov/ProveIt/tree/8c6eec9d19ef80538c8e61ff6e0180040583446a/SetTheory/Cardinals/docs/reports/hilbert-tenth-problem/con-zfc-diophantine-forms}{Pinned repository directory}.
\bibitem{signed17} ProveIt package 29, \emph{Seventeen unrestricted integers:
complete new reduction}, \texttt{signed17/SIGNED17\_PROOF.md}, at \cite{repo}.
\bibitem{quartic31} ProveIt package 30, \emph{A 31-positive-variable quartic
by eliminating a packing register}, \texttt{arithmetic/QUARTIC31\_PROOF.md};
and \emph{Nine-variable equation from the compact coefficients and the
31-variable quartic}, \texttt{nine/COMPACT\_NINE\_PROOF.md}, at \cite{repo}.
\bibitem{positive12} ProveIt package 30, \texttt{positive/POSITIVE12\_PROOF.md}
and its arithmetic circuits, at \cite{repo}.
\bibitem{coding} ProveIt package 27, \texttt{nine/build\_improved.py},
complete-homogeneous and multivariate Horner evaluation, at \cite{repo}.
\bibitem{jones} James P. Jones, \emph{Universal Diophantine equation},
Journal of Symbolic Logic \textbf{47} (1982), no. 3, 549--571.
\href{https://doi.org/10.2307/2273588}{doi:10.2307/2273588}.
\bibitem{sun} Zhi-Wei Sun, \emph{Ten unknowns for Hilbert's tenth problem
over the integers}, arXiv:2609.01594v2 (2026), Theorem 1.1.
\url{https://arxiv.org/abs/2609.01594}.
\bibitem{endpoint} ProveIt contributors,
\path{Computability/BusyBeaver/Lean/BusyBeaver/ZFC/KEndpoint215.lean}
and \path{Computability/HilbertTenthProblem/Lean/Diophantine/Paper1982/},
at the same pinned revision. Selected source files accompany the archive.
\end{thebibliography}
\end{document}
